%% Revisao estilistica em 15/11/2026, 34 trechos. Este capitulo concentrava as
%% formulas repetidas do documento, contadas antes da passagem:
%%   3x "Apesar dos avancos significativos, [...] ainda enfrenta desafios
%%      importantes" -- abertura IDENTICA das analises criticas de 3.1.5,
%%      3.2.5 e 3.3.5;
%%   3x "O periodo 2015-2025 testemunhou" -- abertura de 3.2.1, 3.3.1 e 3.4.1;
%%   3x "A integracao efetiva de conhecimento de dominio especifico";
%%   4x "representa um desafio significativo"; 3x "permanece problematica";
%%   2x "Um avanco particularmente relevante"; 2x "Mais recentemente,
%%      arquiteturas baseadas em Transformers"; 23x "significativ*".
%% Cada abertura passou a entrar pelo gargalo tecnico do proprio eixo (custo
%% embarcado, desalinhamento de resolucao e taxa, acoplamento mapa-controle), e
%% cada fechamento, antes uma "direcao promissora para pesquisas futuras",
%% passou a declarar a posicao desta tese em relacao aquele gargalo -- inclusive
%% quando a posicao e nao atacar o problema. Nenhuma referencia foi trocada,
%% acrescentada ou removida nesta passagem; confira com
%%   git diff -U0 -- cap_revisao_bib/revisao.tex | grep -c '\\cite'
%% Renomeado em 13/11/2026. O capitulo nao executa protocolo de revisao sistematica
%% -- nao ha bases, consultas, datas, criterios de inclusao, numero de registros nem
%% processo de selecao documentados -- e chamar-se sistematico prometia um metodo que
%% nao foi realizado. Passa a declarar o que de fato e: revisao critica, com o critério
%% de selecao explicitado abaixo e com o alcance da selecao reconhecido.
%% Titulo alterado em 15/11/2026. Era "Revisao Critica da Literatura" -- o unico
%% titulo da tese que afirmava uma QUALIDADE em vez de nomear o conteudo, ao lado de
%% "Fundamentacao Teorica", "Avaliacao Experimental", "Resultados e Discussao" e
%% "Conclusao". Criticidade e coisa que o leitor julga. As quatro subsecoes chamadas
%% "Analise Critica e Desafios Futuros" passaram a nomear o que de fato discutem.
\chapter{Revisão da Literatura}
\label{cap:revisao}

A inspeção de subestações elétricas é tarefa predominantemente humana em ambiente energizado, e é desse fato que decorre o interesse por automatizá-la: há limite de segurança para a aproximação a equipamento sob tensão, os eventos de interesse são esparsos ao longo de rotas fixadas por norma, e o custo por inspeção é alto. Os \ac{VTA}, também chamados \ac{RMT} ou \ac{UGV}, são a plataforma em que essa automação tem sido tentada, e a navegação autônoma é a função de que as demais dependem.

Esta revisão analisa criticamente o estado da arte na integração de Aprendizado por Reforço, Visão Computacional e Fusão de Sensores para navegação autônoma de \ac{UGV} em ambientes industriais desafiadores, com ênfase no contexto de subestações elétricas, no período entre 2015 e 2025.

É uma \emph{revisão crítica}, e não uma revisão sistemática: não se executou protocolo prospectivo com bases, consultas, datas de busca, critérios de inclusão e exclusão, número de registros recuperados e processo de seleção documentado. A seleção foi feita por amostragem orientada ao problema: trabalhos recuperados em busca temática nas três áreas e, dentro delas, os que (i) tratam navegação ou inspeção robótica autônoma, (ii) são citáveis por identificador verificável e (iii) sustentam ou contradizem alguma afirmação feita nesta tese. O alcance dessa seleção é, portanto, limitado: ela não autoriza alegação de cobertura exaustiva da literatura, e as lacunas apontadas adiante devem ser lidas como lacunas \emph{nesta amostra}, não como ausências demonstradas no campo.

Os trabalhos citados ao longo do capítulo foram conferidos individualmente por identificador (DOI ou endereço resolvível) quanto a autoria, título, veículo e ano, e as afirmações atribuídas a cada um foram confrontadas com o conteúdo do trabalho.

\section{Aprendizado por Reforço (RL) e Aprendizado por Reforço Profundo (DRL)}

Enquanto o Capítulo 2 apresentou os fundamentos teóricos do Aprendizado por Reforço, esta seção analisa criticamente sua evolução e aplicações específicas em navegação autônoma de \ac{UGV} no período 2015-2025, com foco em ambientes industriais complexos como subestações elétricas.

\subsection{Evolução dos Algoritmos DRL para Navegação Robótica}

A trajetória dos algoritmos de \ac{DRL} para navegação robótica responde a uma limitação precisa dos métodos clássicos, a de exigir um modelo do ambiente que o robô raramente possui. %% Corrigido em 14/11/2026 (auditoria bibliografica, duas correcoes nesta frase):
%% (i) Tai2017 era descrito como uso de DQN. O trabalho e de CONTROLE CONTINUO,
%% treinado por DDPG assincrono -- DQN opera sobre acoes discretas e nao e o
%% algoritmo da fonte; (ii) dizia-se que Zhu2017 "expandiu" Tai2017, e mais adiante
%% nesta mesma revisao dizia-se o inverso. Os dois trabalhos sao de 2017 e
%% independentes; a genealogia nao e demonstravel e foi substituida por comparacao
%% factual entre as abordagens.
\citeonline{Tai2017} demonstraram navegação sem mapa de robôs móveis por aprendizado
profundo com \emph{controle contínuo}, treinada por \ac{DDPG} assíncrono e alimentada
diretamente por leituras esparsas de laser, com transferência da política do simulador
para o robô real. \citeonline{Zhu2017}, no mesmo ano e por via independente, trataram
a navegação dirigida a alvo a partir de imagens, dispensando igualmente mapa explícito,
mas tomando como entrada a cena visual e o alvo desejado em vez de distâncias.

\citeonline{Schulman2017} propuseram o \ac{PPO}, que se tornou padrão na área pela estabilidade do passo de atualização. \citeonline{Haarnoja2018} introduziram o \ac{SAC}, que incorpora um termo de entropia na função objetivo, resultando em políticas mais robustas e explorativas. Análises comparativas realizadas por \citeonline{Ibarz2021} demonstraram que algoritmos baseados em política como \ac{PPO} e \ac{SAC} geralmente superam métodos baseados em valor como \ac{DQN} em tarefas de navegação robótica com espaços de ação contínuos.

\subsection{Abordagens Sim-to-Real para Navegação Autônoma}

Política treinada apenas em simulação degrada-se no robô físico por efeito do \textit{reality gap}, a discrepância entre a dinâmica modelada no simulador e o comportamento dos sensores e atuadores em operação. \citeonline{Tobin2017} introduziram o \textit{Domain Randomization}, que varia deliberadamente os parâmetros da simulação durante o treinamento, tornando a política robusta o suficiente para operar no domínio real.

\citeonline{Peng2018} expandiram esta abordagem com o conceito de \textit{Dynamics Randomization}, focando especificamente na variabilidade dos parâmetros físicos. \citeonline{Sadeghi2017} demonstraram navegação visual sem necessidade de dados reais durante o treinamento, utilizando apenas ambientes simulados com randomização visual extensiva.
%% Frase removida em 13/11/2026 (auditoria bibliografica): atribuia a um trabalho
%% nao localizado em nenhuma base uma abordagem hibrida de randomizacao de dominio
%% com adaptacao adversarial.

Estudos comparativos realizados por \cite{Hwangbo2019} e \citeonline{Loquercio2021} demonstraram que abordagens Sim-to-Real bem projetadas podem atingir desempenho comparável ou superior a métodos treinados diretamente no mundo real, com fração do custo e tempo de desenvolvimento.

%% Ligacao acrescentada em 15/11/2026. Esta subsecao quase foi removida por nao ter
%% ocorrencia nos Capitulos 4 a 7 -- criterio que estava errado aqui: ela e o que
%% posiciona a limitacao de escopo mais importante da tese, declarada na Secao 7.3
%% ("a avaliacao e em simulacao"). O criterio certo nao e se o tema reaparece, e sim
%% se alguma afirmacao do trabalho depende dele.
Esta frente é a que situa o limite de escopo desta tese. A avaliação do
Capítulo~\ref{cap:avaliacao} é inteiramente simulada e nenhuma transferência para
\textit{hardware} foi tentada, de sorte que o que aqui se mede é a lógica de decisão da
arquitetura, e não o seu desempenho em pátio. A literatura acima descreve o que separa
uma coisa da outra, e é contra ela que a Seção~\ref{sec:conc-escopo} delimita o alcance
dos resultados.

\subsection{Aprendizado por Reforço Seguro para Ambientes Industriais}

A segurança é uma preocupação primordial na navegação autônoma em ambientes industriais como subestações elétricas. \citeonline{Garciaetal2015} levantaram os métodos de \ac{RL} Seguro e os organizaram em três famílias, conforme o que cada uma modifica: a função de recompensa, o processo de exploração ou o critério de otimização.

\citeonline{Achiam2017} introduziram o \ac{CPO}, um algoritmo que garante restrições de segurança durante o processo de otimização da política. \citeonline{Chow2018} expandiram esta abordagem com o conceito de \ac{CMDP}, onde o objetivo é maximizar a recompensa enquanto se satisfaz restrições de custo. %% Corrigido em 14/11/2026: a frase atribuia a Thananjeyan2021 funcoes de barreira e
%% garantia de invariancia em conjunto seguro, e situava o trabalho em navegacao
%% industrial. O titulo da propria entrada e "Recovery RL: safe reinforcement learning
%% with LEARNED recovery zones": o mecanismo e um critico de seguranca APRENDIDO que
%% alterna entre politica de tarefa e politica de recuperacao. Nada ali garante
%% invariancia, e os ensaios nao sao de navegacao industrial. Garantia formal nao se
%% atribui a uma fonte sem que a fonte a demonstre.
\citeonline{Thananjeyan2021} seguiram outro caminho: em vez de restringir a
otimização, aprenderam um crítico de segurança a partir de violações observadas e o
usaram para alternar entre uma política de tarefa e uma política de recuperação quando
o risco estimado cresce. A salvaguarda é, portanto, aprendida e probabilística: não
oferece a garantia de invariância que uma função de barreira persegue.

Para ambientes com requisito de segurança explícito, \citeonline{Brunke2022} combinam \ac{RL} com controladores de segurança baseados em modelo que podem intervir quando a política aprendida propõe ações potencialmente perigosas. Esta abordagem híbrida permite explorar os benefícios do aprendizado adaptativo enquanto mantém garantias formais de segurança.

%% Ligacao acrescentada em 15/11/2026, pela mesma razao da subsecao anterior: e o que
%% da sentido a frase da Secao 7.2 -- "esses sao os limites da evidencia de seguranca,
%% e nao uma garantia de seguranca".
A distinção entre essas vias importa ao que esta tese pode e não pode afirmar. A
precedência categórica de segurança do Capítulo~\ref{cap:proposta} é regra de projeto,
verificável por inspeção do código, e não garantia probabilística aprendida nem
invariância demonstrada por função de barreira. O que o
Capítulo~\ref{cap:resultados} apresenta é evidência empírica sobre o comportamento dessa
regra num conjunto de execuções, com os seus intervalos declarados --- e é precisamente
por existirem as garantias formais descritas acima que convém dizer que não se
reivindica nenhuma delas.

\subsection{Aplicações Específicas em Subestações Elétricas}

A aplicação de \ac{DRL} para navegação de \ac{UGV} em subestações elétricas apresenta desafios particulares devido à interferência eletromagnética, obstáculos metálicos e requisitos rigorosos de segurança. O ambiente eletromagnético da instalação é, por si, objeto de estudo e de indexação própria na literatura de compatibilidade eletromagnética de subestações \cite{He2007EMCSubstation}, e o estado da pesquisa e do desenvolvimento em robôs de inspeção para esse ambiente foi recentemente revisado por \citeonline{Wei2024Advancements}.

%% Reescrito em 13/11/2026 (auditoria bibliografica). A referencia anterior deste
%% paragrafo nao foi localizada no Crossref nem no OpenAlex; a segunda existe, mas
%% trata de fusao DGPS/dead reckoning para estimacao de pose, e nao de integracao
%% entre DRL e navegacao classica. As duas frases foram refeitas sobre o que os
%% trabalhos efetivamente publicados dizem.
\citeonline{Ou2024SensorFusionDRL} combinaram \ac{DRL} com fusão de sensores para navegação autônoma de robô móvel, ilustrando a frente em que a política de movimento é aprendida a partir de observações já fundidas. Em subestações, os sistemas descritos em operação seguem via distinta: \citeonline{Yan2024CombinedNav} relatam navegação combinada por fusão DGPS/\textit{dead reckoning} para robô de inspeção, na qual o objeto da fusão é a estimação de pose, e não a política de decisão.

Estudos de caso conduzidos por \cite{Lactec2022} e \cite{Iberdrola2022} em subestações reais demonstraram a viabilidade prática destas abordagens, e relatam ganho de eficiência e de cobertura de inspeção sobre os métodos tradicionais, sem decompor esse ganho por mecanismo.

\subsection{Limites do Aprendizado por Reforço neste Domínio}

A escalabilidade é o limite mais visível desta linha de trabalho. A maioria dos estudos revisados opera em áreas pequenas ou simplificadas, e adaptar a política a uma mudança de \textit{layout} ou de condição operacional exigiria aprendizado continuado, que nenhum dos trabalhos aqui examinados sustenta em operação.

A interpretabilidade das políticas aprendidas é outra limitação crítica, especialmente em contextos industriais onde a compreensão das decisões do sistema é essencial para aceitação e certificação. \citeonline{Juozapaitis2022} exploram a explicação de políticas de \ac{DRL} por decomposição da recompensa, mas a aplicação desse tipo de método em sistemas de navegação completos ainda é limitada.

Para o recorte desta tese, a consequência é de projeto. A opacidade da política aprendida e o custo de retreinamento são as duas razões pelas quais a camada de supervisão do Capítulo~\ref{cap:proposta} arbitra por limiar determinístico, e pelas quais a tabela \textit{Q} comparece apenas como linha de base aprendida (Seção~\ref{sec:setup-qdriven}).

%% Figura removida em 13/11/2026. Era infografico gerado em ferramenta externa, com
%% marca d'agua "Made with Napkin" impressa, trofeu decorativo e rotulos em ingles
%% ("Pros"/"Cons"). Alem da forma: (i) listava dez propriedades do DRL sem citacao
%% alguma, com fonte "O proprio autor", em capitulo de revisao critica -- o mesmo
%% padrao de afirmacao sem lastro corrigido no restante do capitulo; (ii) trazia
%% "Transferencia Sim-to-Real" como PRO e "Desafios de transferencia" como CONTRA,
%% ambos definidos como transferir politicas de simulacao para o mundo real;
%% (iii) nao correspondia ao texto, que elenca interpretabilidade e nao "desafios de
%% recompensa". Nao era referenciada em lugar nenhum. O arquivo permanece em figuras/.

%% Secao COMPRIMIDA em 15/11/2026, de 832 para cerca de 300 palavras. Era o eixo mais
%% pesado do capitulo (19%) e a tese usa visao UMA vez nos Capitulos 4 a 7. Mantido o
%% que tem consequencia: o reconhecimento de anomalia, porque os eventos desta tese SAO
%% anomalias, e o obstaculo de dominio. Removido o bloco de navegacao e localizacao
%% visual, que a Secao 2.5.4 agora cobre pelo lado da localizacao em mapa dado, e
%% reduzida a linhagem de arquiteturas de rede, que nao reaparece em lugar nenhum.
\section{Visão Computacional para Navegação Autônoma}

A visão computacional entrega a um \ac{UGV} o tipo de informação que o \ac{LiDAR} não
fornece, a identidade semântica do que está à frente, e não apenas a sua geometria. A
linhagem de arquiteturas que a sustenta foi reorganizada pelo custo computacional
embarcado, das \ac{CNN} residuais \cite{He2016} aos \textit{Transformers} de visão
\cite{Carion2020, Liu2021SwinTransformer}, que superam as primeiras em \textit{benchmarks}
de percepção ao preço de mais computação.

Para a inspeção, porém, o problema decisivo não é chegar ao equipamento, e sim reconhecer
a anomalia uma vez diante dele. Essa frente concentra o trabalho mais próximo desta tese,
ainda que por contraste: detecção de defeito em superfície industrial
\cite{Ullah2020DefectDetection}, detecção de anomalia com pouca ou nenhuma rotulagem
\cite{Bergmann2020}, fusão de imagem infravermelha e visível para equipamentos de
subestação \cite{Chi2024LMDFusion} e detectores leves para alvos de pequena extensão, nos
quais a taxa de detecções perdidas é alta \cite{Wang2024SmallTargets}.

Dois obstáculos separam o resultado de \textit{benchmark} do desempenho em pátio. O
primeiro é a degradação sob reflexo metálico e variação extrema de iluminação, condições
correntes numa subestação a céu aberto e às quais a segmentação semântica é sensível
\cite{Pfeuffer2019, Cai2025AttentionSeg}. O segundo é estrutural: a anomalia que se quer
detectar é rara no equipamento que se quer inspecionar, de sorte que dados rotulados de
falha escasseiam por definição.

Esta tese não ataca esse gargalo. A percepção é estipulada, com os eventos de interesse
entrando no sistema como leituras já classificadas (Seção~\ref{sec:prop-sensoriamento}),
e o objeto de estudo é o que o robô faz com elas.
%% Secao COMPRIMIDA em 15/11/2026, de 589 para cerca de 260 palavras. Esta tese nao
%% funde modalidade alguma: a arquitetura consome evento ja rotulado. As quatro
%% subsecoes viraram prosa corrida, e o que se preservou foi o que tem consequencia
%% adiante -- a distincao entre fundir ANTES e DEPOIS da decisao, que e o que situa o
%% objeto desta tese, e a dessincronizacao temporal entre modalidades de taxa
%% distinta, que e o vizinho direto do defeito de relogio da Secao 6.13.
\section{Fusão Sensorial para Navegação Robusta}

A fusão sensorial existe porque nenhuma modalidade isolada cobre as condições de uma
subestação, e a literatura recente organiza-se em torno do desalinhamento entre
resolução, campo de visão e taxa de amostragem das modalidades que se quer combinar
\cite{Chen2017, Jaritz2019}. A atenção cruzada deslocou o problema da escolha de um ponto
de fusão para a aprendizagem de \emph{onde} fundir \cite{Prakash2021}. Em instalação
energizada acrescenta-se a interferência eletromagnética, que degrada cada modalidade de
maneira distinta e exige ponderação sensível à confiabilidade momentânea de cada leitura
\cite{Yan2024CombinedNav, Alatise2020Review}.

Convém separar dois pontos em que a fusão entra na inspeção autônoma, porque só o segundo
toca esta tese. Antes da decisão, ela compõe a estimativa de posição, e é aí que se
concentra a literatura de navegação citada acima. Depois dela, compõe o \emph{diagnóstico}
do equipamento: para descargas parciais, por exemplo, métodos elétricos e ópticos são
combinados com detecção por rede neural \cite{Monzon2025PartialDischarge}, porque nenhuma
das vias localiza a descarga isoladamente.

Três limites atravessam essa literatura. A calibração extrínseca deriva em operação, com
temperatura e vibração alterando parâmetros estimados em bancada. A dessincronização
temporal entre modalidades de taxa distinta é o segundo, e o mais próximo do que este
trabalho encontrou na prática: compor leituras colhidas em bases de tempo que não
coincidem produz erro que nenhuma ponderação corrige, e a Seção~\ref{sec:res-verificacao-relogio}
mede um caso dessa família. O terceiro, a escalabilidade em número de modalidades, pesa
menos aqui, porque os sistemas revisados param em três ou quatro sensores.

Nenhuma dessas técnicas é implementada nesta tese. A arquitetura do
Capítulo~\ref{cap:proposta} recebe o evento já formado, com modalidade e intensidade
declaradas, e o que se estuda é a decisão que vem depois.
\section{Integração de Técnicas para Sistemas Completos}
\label{sec:rev-integracao}

Integrar aprendizado, visão e fusão num mesmo sistema introduz um problema que nenhuma das três áreas resolve isoladamente, o de atribuir a cada componente a parte do desempenho que lhe cabe. Esta seção analisa criticamente as abordagens de integração propostas na literatura recente, com foco em aplicações para subestações elétricas.

\subsection{Arquiteturas Integradas para Navegação Autônoma}

O acoplamento entre representação do mapa e controle hierárquico é o que distingue as arquiteturas integradas entre si. %% Corrigido em 14/11/2026: este trecho afirmava que Tai2017 expandiu Zhu2017,
%% enquanto a secao anterior afirmava o contrario. Ambos sao de 2017; nenhum expande
%% o outro. Mantem-se a comparacao de entrada sensorial, que e verificavel.
\citeonline{Zhu2017} propuseram uma das primeiras arquiteturas que combinava visão
computacional com aprendizado por reforço profundo para navegação dirigida a alvos
visuais. \citeonline{Tai2017}, de forma contemporânea e independente, partiram de
leituras esparsas de laser em vez de imagens, com o que trocaram riqueza semântica da
entrada por um espaço de observação de dimensão muito menor.

%% Reescrito em 13/11/2026: o mapeamento do Active Neural SLAM e metrico aprendido,
%% nao semantico, e o Occupancy Anticipation nao incorpora memoria episodica -- ele
%% antecipa a ocupacao de regioes ainda nao observadas.
\citeonline{Chaplot2020} introduziram uma arquitetura hierárquica --- política global, planejador analítico e política local --- sobre um mapeamento métrico aprendido, obtendo exploração e navegação eficientes em ambientes previamente inexplorados. \citeonline{Ramakrishnan2020Occupancy} atacaram o mesmo problema pelo lado da representação, antecipando a ocupação de regiões ainda não observadas a partir do que está no campo de visão, o que reduz o custo de exploração necessário para navegar.

Para aplicações em ambientes industriais, \citeonline{Kahn2021} propuseram uma arquitetura que integra percepção visual, fusão sensorial e aprendizado por reforço com garantias de segurança, demonstrando navegação robusta mesmo em presença de obstáculos dinâmicos e condições ambientais variáveis.
%% Corrigido em 13/11/2026: Yan et al. (2024) nao empregam aprendizado por reforco.
%% O trabalho e um sistema de navegacao combinada DGPS/DR por fusao multissensorial
%% para robo de inspecao de subestacao -- posicionamento, nao decisao.
Para subestações especificamente, o que a literatura revisada reporta é sobretudo
integração voltada ao \emph{posicionamento}: \citeonline{Yan2024CombinedNav} combinam
DGPS e odometria por \textit{dead reckoning} em um sistema de navegação para robô de
inspeção de subestação, e \citeonline{Fan2024MultiSensorSubstation} aplicam fusão
multissensorial ao posicionamento preciso do mesmo tipo de plataforma. Nenhum dos dois
trata a decisão de \emph{deixar a rota} em resposta ao que se percebe, que é o objeto
desta tese.

\subsection{Frameworks Modulares e Adaptáveis}
\label{sec:rev-modulares}

Os \textit{frameworks} modulares respondem à necessidade de trocar um componente sem reconstruir o sistema. \citeonline{Savva2019} introduziram o Habitat, uma plataforma para pesquisa em navegação visual incorporada que facilita a integração de diferentes algoritmos de percepção e controle. \citeonline{Wijmans2020} utilizaram esta plataforma para demonstrar navegação ponto-a-ponto em ambientes complexos utilizando apenas sensores visuais.

%% Corrigido em 14/11/2026: o texto apresentava o RoboTHOR (2020) como expansao do
%% ManipulaTHOR (2021), sequencia cronologicamente impossivel. O ano de Deitke2020
%% havia sido corrigido antes, mas a relacao temporal permaneceu errada.
%% Os dois sao extensoes irmas da mesma plataforma, em eixos distintos.
\citeonline{Deitke2020RoboTHOR} introduziram o RoboTHOR, que acrescenta à plataforma
um par entre cenas simuladas e suas contrapartes físicas, voltado a medir a perda de
desempenho na transferência do simulador para o robô real.
%% Acrescentado em 15/11/2026. Estas plataformas sao a resposta da area ao problema de
%% denominador comum levantado na Secao 3.5, e omiti-las tornaria aquela secao injusta:
%% o denominador existe, mas para outro objeto. Removido no mesmo passo o ManipulaTHOR,
%% que estende a plataforma para manipulacao e nao tem relacao com este trabalho.
Essas plataformas são a resposta mais direta ao problema de comparabilidade retomado na
Seção~\ref{sec:rev-reporte}: oferecem cena, tarefa e métrica fixas, sobre as quais
sistemas distintos podem ser confrontados. O denominador comum existe, portanto --- mas
para a navegação visual guiada por objetivo, que é outro objeto. Não há equivalente para
a camada que decide \emph{quando} interromper uma tarefa em favor de outra, e é por isso
que o Capítulo~\ref{cap:avaliacao} precisa construir o seu.

%% Reescrito em 13/11/2026. A redacao anterior descrevia a Iberdrola2022 como um
%% "framework modular com configuracao flexivel de sensores" e a Lactec2022 como uma
%% extensao dele com conhecimento de dominio. Nenhuma das duas fontes afirma isso:
%% sao, respectivamente, um comunicado de imprensa sobre o robo Antecursor I e uma
%% nota institucional sobre o desenvolvimento de um robo de inspecao. A cadeia de
%% "expandiram esta abordagem" entre elas nao existe na literatura.
No lado industrial, o que se encontra documentado são relatos institucionais de
plataformas em operação, e não arquiteturas descritas em detalhe metodológico: a
Iberdrola divulga o robô Antecursor~I, que transmite imagens de inspeção em tempo real
a partir de subestações \cite{Iberdrola2022}, e o Lactec noticia o desenvolvimento de
um robô para inspeção em subestações de extra alta tensão \cite{Lactec2022}. Esses
relatos atestam o interesse operacional pelo problema, mas são comunicações
corporativas: não informam a arquitetura de decisão, os comparadores empregados nem as
métricas obtidas, de modo que não podem ser tomados como evidência comparável à da
literatura revisada por pares e são citados aqui apenas como indicação de demanda
prática.

\subsection{Integração com Sistemas de Gerenciamento de Ativos}

%% Reescrito em 13/11/2026. A versao anterior atribuia sistemas e arquiteturas a tres
%% fontes que nao os descrevem -- um blog de fornecedor sobre automacao de subestacoes
%% digitais, uma pagina de produto e uma chamada de desafio de P&D -- e encadeava
%% entre elas um "expandiram esta abordagem" inexistente.
A integração entre inspeção robótica e plataformas de gerenciamento de ativos é
apontada como caminho natural de aplicação, mas a documentação disponível é
predominantemente industrial, e não científica. A automação da subestação digital e a
padronização de comunicação que a sustenta --- da qual a IEC~61850 é o exemplo
corrente --- são descritas em material de fornecedores e de integradores
\cite{Clouglobal2024,Selinc2024}, enquanto a inspeção remota e autônoma em subestação
de alta tensão com robôs terrestres aparece formulada como \emph{desafio} de pesquisa
e desenvolvimento por operador do setor \cite{Eletrobras2024}. Registra-se, portanto,
a existência da demanda e do substrato de integração, sem atribuir a essas fontes
arquiteturas, resultados ou métricas que elas não reportam. Esta tese não trata dessa
camada de integração: seu objeto é a lógica de decisão que leva o sensor até o evento
e o afasta do risco.

\subsection{Avaliação Comparativa de Sistemas Integrados}

A comparação entre sistemas integrados é dificultada pela diversidade de ambientes, tarefas e métricas, que raramente coincidem entre dois trabalhos. \citeonline{Anderson2018} propuseram o benchmark Room-to-Room (R2R) para navegação visual baseada em linguagem natural, estabelecendo um padrão para comparação de diferentes abordagens.

\citeonline{Savva2019} expandiram esta ideia com o PointNav, um benchmark para navegação ponto-a-ponto em ambientes desconhecidos utilizando apenas sensores visuais e inerciais. \citeonline{Kadian2020} introduziram métricas específicas para avaliar a transferibilidade de políticas de navegação da simulação para o mundo real, um aspecto crítico para aplicações práticas.

%% Reescrito em 13/11/2026. A versao anterior apresentava um relatorio de mercado
%% como "framework de avaliacao" e a revisao de Wei e Rey como uma extensao dele.
%% Nenhuma das duas fontes propoe metricas de avaliacao.
Para inspeção de subestações, em contraste, \emph{não} se identificou nesta amostra um
\textit{benchmark} consolidado: não há conjunto de cenários, métricas e linhas de base
compartilhados que permita comparar dois supervisores de inspeção publicados por grupos
distintos. O que existe são um levantamento de mercado do segmento de robôs de inspeção
de subestações \cite{GlobalGrowthInsights2024}, útil como indicação de adoção e não
como instrumento de avaliação, e a revisão de \citeonline{Wei2024Advancements} sobre o
estado da pesquisa e do desenvolvimento desses robôs. Essa ausência tem consequência
direta sobre o desenho experimental desta tese: na falta de um \textit{benchmark}
externo, a comparação é construída \emph{internamente}, por ablações sobre um mesmo
controlador compartilhado (Capítulo~\ref{cap:avaliacao}), o que torna cada contraste
interpretável sem depender de uma referência comum que o campo ainda não possui.

\subsection{Lacunas Comuns às Três Frentes}

Falta a esta literatura um ponto de comparação comum. Cada trabalho integra aprendizado, visão e fusão à sua maneira, relata ganho agregado sobre uma linha de base própria e não decompõe esse ganho por componente, de modo que os resultados não são comparáveis entre si nem atribuíveis internamente.

A escalabilidade para complexos de subestação permanece limitada pela área de ensaio, e a ausência de \textit{benchmark} padronizado impede dizer se o limite observado é do método ou do ambiente em que foi medido.

A resposta a falha sensorial e a evento imprevisto é afirmada com mais frequência do que medida. Poucos dos sistemas revisados relatam o que ocorre quando uma modalidade cai ou quando um obstáculo não previsto entra na rota, e quase nenhum separa a contribuição do mecanismo de detecção da do mecanismo de recuperação.

%% Reescrito em 15/11/2026. O paragrafo anterior era o unico da secao que nao nomeava
%% lacuna alguma: fechava em "sera essencial para a adocao ampla", que e previsao sobre
%% a industria e nao achado sobre a literatura, e era tambem o unico sem citacao. A
%% versao atual nomeia a exigencia que a avaliacao por desempenho medio nao atende e
%% liga essa restricao a escolha de projeto desta tese.
Por fim, a certificação operacional em instalação energizada impõe uma exigência que a avaliação por desempenho médio não atende: declarar limite de garantia e condição de falha. Política aprendida fim a fim oferece pouca transparência sobre quando a reação de segurança cede lugar à tarefa, e sem protocolo de verificação formal não se demonstra conformidade normativa. É dessa restrição que decorre a opção desta tese por uma regra de precedência categórica de segurança explícita, cujo efeito e cujo custo podem ser auditados em separado.

%% Secao criada em 13/11/2026. A revisao discutia as tres areas em paralelo, mas em
%% nenhum ponto confrontava esta tese com os trabalhos que estao de fato mais perto
%% dela -- supervisores de missao e arquiteturas hibridas com arbitragem. Sem esse
%% confronto, a contribuicao so podia ser lida como soma de tecnologias conhecidas.
%% Secao NOVA em 15/11/2026, e e a que faltava. A contribuicao declarada no
%% Capitulo 7 e METODOLOGICA -- a decomposicao por ablacao com analise que declara as
%% proprias fronteiras --, mas o capitulo nao revisava pratica metodologica alguma:
%% "tamanho de efeito", "Holm" e "censura" ocorriam ZERO vezes aqui contra 14, 24 e 17
%% nos capitulos posteriores. A afirmacao de que esta literatura "nao decompoe o ganho
%% por componente" estava ASSERIDA e nao demonstrada, de sorte que a pedra sobre a qual
%% a contribuicao se apoia era a unica do documento sem apoio proprio.
%% O argumento foi construido sobre um fato verificavel NESTE capitulo, e nao sobre
%% alegacao a respeito do que cada artigo reporta, que exigiria auditoria que nao foi
%% feita -- e a limitacao esta declarada no ultimo paragrafo.
\section{O Que Esta Literatura Reporta}
\label{sec:rev-reporte}

As quatro frentes revisadas acima resumem $66$ trabalhos, e vale observar o que esses
resumos \emph{não} contêm: nenhuma grandeza comparável. Não há nestas páginas um único
tamanho de efeito, intervalo de confiança, número de execuções por condição ou ganho
percentual atribuído a um componente isolado. A ausência é parcialmente uma escolha
deste texto, que privilegiou o mecanismo sobre o número; mas é também a razão de a
escolha ter sido possível sem perda, e é essa segunda parte que importa.

A dificuldade é de denominador. Cada trabalho define a sua própria plataforma, a sua
própria linha de base e o seu próprio cenário, de modo que o ganho que relata está
expresso em unidades que só valem dentro dele. Um sistema que melhora a taxa de detecção
em relação a uma versão anterior de si mesmo, outro que reduz o tempo de percurso em
relação a um operador humano e um terceiro que eleva a acurácia de classificação sobre um
conjunto próprio não produzem, juntos, nenhuma ordenação. Transcrever esses números lado
a lado daria a aparência de comparação sem a sua substância.

Disso decorre a consequência que esta tese enfrenta. Quando o sistema avaliado integra
várias técnicas --- e as arquiteturas da Seção~\ref{sec:rev-integracao} integram, por
construção ---, um ganho relatado no agregado não diz qual das partes o produziu. A
questão não é se o conjunto funciona, e sim quanto cada peça contribui, porque é essa a
informação de que precisa quem for reusar uma delas. Responder a isso exige um
denominador comum construído de propósito: o mesmo controlador, o mesmo ambiente, o mesmo
protocolo de eventos, e variação de um mecanismo por vez. É o desenho adotado no
Capítulo~\ref{cap:avaliacao}, e a razão de ele existir está aqui.

A segunda consequência é sobre o que acompanha o número quando ele existe. Atribuir
efeito a mecanismo em simulação estocástica exige declarar a incerteza da atribuição ---
quantas execuções a sustentam, que fração delas foi interrompida sem concluir, e quanto
da diferença observada sobrevive à correção para o número de comparações feitas. Essas
declarações são o que permite distinguir um efeito de uma flutuação, e é por isso que o
Capítulo~\ref{cap:resultados} as apresenta para cada contraste, com as suas próprias
fronteiras ditas em voz alta.

%% Ressalva acrescentada em 15/11/2026, logo depois de escrita a secao: a afirmacao de
%% que "nao ha base comum" era forte demais e a propria Secao 3.4.2 a contradizia, ao
%% descrever Habitat e RoboTHOR, que sao exatamente tentativas de base comum. A
%% afirmacao correta e mais estreita, e continua suficiente.
Duas ressalvas delimitam o que se afirmou. A primeira é que base comum \emph{existe} em
parte do campo: as plataformas de referência da Seção~\ref{sec:rev-modulares} fixam cena,
tarefa e métrica, e sobre elas sistemas distintos são de fato confrontáveis. Elas cobrem,
porém, a navegação visual guiada por objetivo, e não a camada que decide quando
interromper uma tarefa --- que é o objeto desta tese, e para o qual não se localizou
equivalente.

A segunda é que não se conduziu auditoria sistemática das práticas estatísticas das
fontes primárias, e nada do que se disse acima afirma que elas reportem mal dentro dos
seus próprios termos. O que se afirma é mais estreito: para o objeto aqui estudado, a
literatura revisada não oferece base comum sobre a qual os ganhos relatados possam ser
confrontados, e por isso a atribuição por componente precisa ser produzida, e não
consultada.

\section{Trabalhos Mais Próximos}
\label{sec:rev-proximos}

As três áreas revisadas acima situam o contexto, mas os trabalhos \emph{mais próximos}
desta tese não são os que percebem melhor nem os que aprendem políticas: são os que
decidem \emph{quando deixar de prosseguir e reagir}. Essa é a função de um supervisor de
missão, e ela tem uma linhagem própria --- das arquiteturas comportamentais e híbridas
clássicas às arquiteturas aprendidas hierárquicas ---, com a qual a proposta deste
trabalho precisa ser confrontada diretamente.

A \autoref{tab:rev-proximos} organiza esse confronto por quatro atributos que são
exatamente os que esta tese manipula experimentalmente: como o conflito entre reagir e
prosseguir é resolvido; se o sinal que motivou a reação é retido em um
\emph{acumulador com decaimento}, e não apenas se há alguma forma de estado
temporal, que quase toda arquitetura tem; se uma camada deliberativa madura é reaproveitada em vez de substituída; e se o
efeito de cada mecanismo de supervisão é atribuído experimentalmente, em vez de
reportado no agregado.

\begin{table}[htb]
\centering
%% Reescrito em 14/11/2026. A coluna chamava-se "persistencia temporal
%% explicita" e marcava "Nao" para Brooks1986 -- mas a subsuncao usa maquinas de
%% estado AUMENTADAS, com estado interno e constantes de tempo na supressao e na
%% inibicao. Ausencia de acumulador com decaimento nao e ausencia de persistencia
%% temporal. A coluna passa a nomear o mecanismo especifico que esta tese manipula, e
%% a legenda distingue os quatro estatutos possiveis de uma celula.
\caption{Trabalhos mais próximos, por atributo de supervisão. A terceira coluna
pergunta por um mecanismo \emph{específico} --- o sinal que motivou a reação ser retido
em um acumulador que decai ---, e não por persistência temporal em geral: estado
interno, temporizadores e mapas acumulados são formas de persistência que vários destes
trabalhos têm. As células distinguem quatro estatutos: \emph{não possui} o mecanismo;
\emph{não o documenta} como mecanismo; \emph{usa outro mecanismo}, que é então
nomeado; e \emph{não aplicável} ao desenho. A última coluna distingue reportar o
desempenho do sistema completo de atribuir efeito a um mecanismo de supervisão isolado
--- nos trabalhos de política aprendida, o objeto do estudo é a política, e a
arbitragem entre reagir e prosseguir não é um mecanismo separável sobre o qual se possa
ablacionar. Trabalhos agrupados em uma mesma linha compartilham o estatuto nestes
quatro atributos, e não são por isso equivalentes em todos os aspectos.}
\label{tab:rev-proximos}
%% 14/11/2026: celulas comprimidas -- "Temporizadores" e "Sim" encostavam.
%% Aumentado o espacamento entre colunas e redistribuida a largura, dando mais espaco
%% a coluna do mecanismo de persistencia, que agora nomeia mecanismos em vez de dizer
%% "Nao". Soma das larguras + espacamento permanece abaixo de \textwidth.
\scriptsize
\setlength{\tabcolsep}{4.5pt}
\renewcommand{\arraystretch}{1.15}
\begin{tabular}{p{0.165\textwidth}p{0.205\textwidth}p{0.185\textwidth}p{0.115\textwidth}p{0.155\textwidth}}
\hline
\textbf{Trabalho / classe} &
\textbf{Conflito entre reagir e prosseguir} &
\textbf{Acumulador do sinal, com decaimento} &
\textbf{Camada deliberativa reaproveitada} &
\textbf{Atribuição entre mecanismos}\\
\hline
\citeonline{Brooks1986} --- controle em camadas &
Supressão e inibição por prioridade fixa de camada &
Outro mecanismo: estado interno em máquinas de estado aumentadas, com constantes de
tempo na supressão e na inibição &
Não (sem planejador) &
Demonstração do sistema\\
\hline
\citeonline{Arkin1989,Arkin1998} --- esquemas motores &
Soma vetorial de comportamentos ponderados por ganhos &
Não possui: a soma é função das percepções do instante &
Não (comportamental) &
Demonstração do sistema\\
\hline
\citeonline{Barraquand1992}; \citeonline{Koren1991} --- campos potenciais &
Gradiente do campo somado; limitações do método analisadas &
Não possui: o campo é função da configuração corrente &
Sim, no caso numérico &
Análise do método\\
\hline
\citeonline{Bonasso1997} --- arquitetura em três camadas &
Sequenciador reativo interpõe-se entre planejador e habilidades &
Temporizadores do sequenciador &
Sim &
Demonstração do sistema\\
\hline
\citeonline{Konolige1997} --- Saphira &
Arbitragem difusa entre comportamentos &
Não documentada como mecanismo &
Sim &
Demonstração do sistema\\
\hline
\citeonline{Zhu2017,Tai2017}; \citeonline{Kahn2021} --- políticas aprendidas fim a fim &
Implícito na política; sem árbitro separável &
Implícita no estado ou no modelo &
Não &
Não é objeto; comparação agregada\\
\hline
\citeonline{Chaplot2020}; \citeonline{Ramakrishnan2020Occupancy} --- hierárquicas aprendidas &
Política global escolhe a meta; planejador analítico a executa &
Mapa acumulado &
Sim (planejador) &
Comparação agregada em \textit{benchmark}\\
\hline
\citeonline{Liu2016VisualNavSubstation}; \citeonline{Yan2024CombinedNav}; \citeonline{Fan2024MultiSensorSubstation} --- inspeção de subestação &
Não tratado: rota fixa; o esforço está no posicionamento &
Não aplicável &
Sim &
Demonstração em plataforma\\
\hline
\textbf{Esta tese} &
Precedência \emph{categórica} de segurança, e não disputa de magnitude &
Memória de canal com limiar e decaimento &
Sim: Nav2/MPPI permanece o único produtor de \texttt{/cmd\_vel} &
Sim, por ablação sobre controlador compartilhado\\
\hline
\end{tabular}
\renewcommand{\arraystretch}{1.0}
\end{table}

Três leituras decorrem da tabela. Primeira: a precedência categórica não é novidade
isolada: a prioridade fixa entre camadas está em \citeonline{Brooks1986} e o
sequenciador de \citeonline{Bonasso1997} cumpre papel análogo. O que não se encontra
nessa linhagem é a \emph{medida do preço} dessa prioridade: quantas inspeções legítimas
ela recusa e quanto tempo de missão ela consome. Segunda: a persistência temporal
explícita do sinal de perigo, por limiar e decaimento sobre um canal dedicado,
distingue a proposta tanto da soma vetorial instantânea dos campos potenciais quanto da
arbitragem comportamental sem memória declarada, e é precisamente o fator que o
contraste contra a heurística genérica isola. Terceira: em nenhuma das classes acima a
atribuição de efeito \emph{entre} mecanismos de supervisão é o objeto do estudo; é o
que esta tese toma como objeto, com o custo de abstrair a percepção para conseguir
fazê-lo.

\section{Conclusão e Implicações para o Framework Proposto}

Esta revisão analisou criticamente o estado da arte na integração de Aprendizado por Reforço, Visão Computacional e Fusão Sensorial para navegação autônoma de \ac{UGV} em ambientes industriais complexos, com foco particular em subestações elétricas. A análise registra progresso em cada área isolada e em sua integração, e identifica limites persistentes que motivam o desenvolvimento de novas abordagens.

%% Reescrito em 13/11/2026. A versao anterior atribuia ao trabalho proposto tres
%% diferenciais -- modularidade configuravel, incorporacao de conhecimento de dominio
%% e mecanismos de deteccao e mitigacao de interferencia eletromagnetica -- que o
%% sistema avaliado nao implementa: ele recebe eventos-oraculo. A delimitacao da
%% Secao~\ref{sec:intro-escopo} passa a valer tambem aqui, com a mesma particao
%% entregue / abstraido / futuro.
Os avanços em algoritmos de \ac{DRL}, combinados com técnicas de transferência
Sim-to-Real, oferecem um paradigma para navegação adaptativa em ambientes dinâmicos; as
arquiteturas modernas de visão computacional ampliam a percepção sob condições
desafiadoras; e as técnicas de fusão sensorial integram modalidades heterogêneas. Esses
avanços formam o contexto em que esta tese se situa, e não o que ela entrega.

O que esta amostra da literatura \emph{não} estabelece, e que motiva o trabalho dos
capítulos seguintes, é \emph{qual mecanismo} de uma arquitetura integrada responde por
qual ganho. Os trabalhos revisados relatam desempenho agregado de sistemas que
percebem, decidem e se desviam simultaneamente; a comparação predominante é contra uma
ronda fixa, que não dispõe de mecanismo para deixar a rota e por isso perde por
construção. É essa lacuna de \emph{atribuição} que a arquitetura de duas camadas do
Capítulo~\ref{cap:proposta} ataca, e é por isso que o Capítulo~\ref{cap:resultados}
desmonta experimentalmente a própria proposta em vez de apenas vencer uma linha de
base.

Em relação ao que a revisão levantou, a delimitação desta tese é a seguinte, nos mesmos
termos da Seção~\ref{sec:intro-escopo}:

\begin{enumerate}
    \item \textbf{Entregue e avaliado}: a supervisão de missão em duas camadas --- a
    regra de prioridade com precedência categórica de segurança, a memória temporal do
    canal de risco com limiar e decaimento, a geração de metas de inspeção e de recuo e
    a integração com a navegação deliberativa ---, avaliada em simulação com ablações
    sobre um controlador compartilhado e atribuição estatística de efeitos a mecanismos.

    \item \textbf{Abstraído}: a detecção, a classificação, a intensidade e a posição
    dos eventos, que entram na arquitetura como eventos-oráculo; e a transformada de
    localização, mantida fixa. São controles de confundimento declarados, adotados para
    que o efeito de cada mecanismo possa ser atribuído, e não capacidades alegadas.

    \item \textbf{Futuro}: a percepção térmica, acústica e visual real alimentando a
    camada de eventos, com tratamento explícito de falsos positivos e negativos; a
    robustez a interferência eletromagnética, que esta tese não implementa nem mede; a
    integração sensorial aprendida; e a validação operacional em instalação real
    (Seção~\ref{sec:conc-futuro}).
\end{enumerate}

Nenhum mecanismo de percepção, de fusão sensorial ou de mitigação de interferência é
alegado como contribuição desta tese. O que se alega é a decomposição experimental de
uma camada de supervisão de inspeção, com os limites de evidência declarados.
